\section{结论与局限性}

本文围绕“A 股分钟级市场状态信息能否预警未来 5 分钟和 10 分钟短时流动性压力”展开研究。综合上述证据，本文发现分钟级市场状态信息能够为未来短时流动性压力提供可验证的样本外预警。RGARCH-CARR-SK 结果说明 MarketLSI 压力风险存在阶段性聚集，并且其估计路径对 realized pressure measure 口径具有敏感性；动态峰度在本样本中主要停留在模型基准附近，因此不作为核心经验发现。QVAR 结果显示尾部条件下市场压力、波动和成交承接变量之间的响应路径不同于中位数状态；SMARTBoost 结果则进一步表明，无泄漏的高频历史特征能够在样本外高风险排序中集中识别未来压力事件。

本文仍存在局限。首先，本文基于 1 分钟 OHLCV 与成交额数据构造短时流动性压力代理指标，重点刻画市场成交、价格振幅、局部波动和成交承接状态共同反映的短时压力变化。其次，QVAR 压力测试是标准化情景模拟，不构成严格因果识别。最后，SMARTBoost 结果体现的是特定时间窗口的样本外预警能力，后续仍可在更长样本、更多市场状态和更高频数据上进一步检验。尽管如此，本文形成了从高频压力指数到动态风险刻画、尾部分位传导和样本外智能预警的完整经验框架。
